% ============================================================
% Amazon-VT CFP 2026-2027 — 1-Page Abstract v2
% Provenance-Aware Agentic Knowledge Graph Systems
% SOA-grounded framing with 2024-2026 citations
% Format: IEEE 2-column
% ============================================================

\documentclass[conference, 10pt]{IEEEtran}

\usepackage[utf8]{inputenc}
\usepackage[T1]{fontenc}
\usepackage{graphicx}
\usepackage{hyperref}
\usepackage{microtype}
\usepackage{cite}
\usepackage{enumitem}
\setlist{nosep, leftmargin=1.5em, topsep=0.2em}

\title{Provenance-Aware Agentic Knowledge Graph Systems\\for Reliable Multi-Source Grounded AI}
\author{}
\date{}

% ============================================================
\begin{document}

\maketitle
\thispagestyle{empty}

\vspace{-1.5em}

\noindent\textbf{Primary Topic:} Knowledge Grounding \hfill
\textbf{Secondary Topic:} Agentic AI

\vspace{0.5em}

Retrieval-augmented generation (RAG)~\cite{Lewis2020RAG} has become a foundational technique for grounding large language model (LLM) inference in external knowledge. Production systems have matured substantially: hybrid retrieval combining semantic (dense) and keyword (sparse) search is now standard practice~\cite{AWSBedrockHybridSearch2024}, and retrieval is increasingly exposed as a tool within agentic workflows where models decide \textit{when} and \textit{what} to retrieve~\cite{OpenAIFunctionCalling2025}. Yet despite these advances, two critical gaps remain. First, even state-of-the-art retrieval systems lack structured reasoning substrates capable of multi-hop synthesis across heterogeneous knowledge sources---a gap that graph-based approaches like GraphRAG~\cite{Edge2024GraphRAG} begin to address but do not fully resolve. Second, current systems provide limited provenance: they can cite documents, but cannot produce fine-grained, verifiable evidence chains linking individual model outputs to specific supporting knowledge, and they lack validation mechanisms that detect and recover from retrieval failures~\cite{Yan2024CRAG,Ji2023HallucinationSurvey}.

This project proposes \textbf{Provenance-Aware Agentic Knowledge Graph Systems (PA-AKG)}, an architecture that unifies structured knowledge graph (KG) reasoning, adaptive multi-source retrieval, and agent-driven validation---capabilities that existing systems provide only in isolation~\cite{Edge2024GraphRAG,Yan2024CRAG,Asai2024SelfRAG}. The core technical contribution is a \textit{provenance graph}: a structured representation maintaining entity- and span-level evidence links between agent outputs and the specific KG triples, document passages, and database records that support them. Unlike document-level citations, provenance graphs enable fine-grained verification---for a multi-hop claim, the graph records which entities, relationships, and source spans contributed, with confidence at each link. Agents dynamically route among retrieval pathways based on task context and provenance completeness. Validation agents assess evidence sufficiency using provenance signals aligned with observability standards~\cite{OpenTelemetryGenAI2024}, triggering corrective retrieval when grounding is weak or contradictory---extending retrieval-control mechanisms~\cite{Yan2024CRAG,Asai2024SelfRAG} from unstructured documents to structured KG reasoning.

\begin{figure}[t]
    \centering
    \includegraphics[width=\columnwidth]{images/architecture-diagram.png}
    \caption{\textbf{PA-AKG Architecture.} Heterogeneous knowledge sources feed a hybrid retrieval layer (vector search + KG queries). An agent orchestration layer coordinates planning, retrieval, and validation agents. A provenance engine records evidence links aligned with observability standards~\cite{OpenTelemetryGenAI2024} for traceable verification.}
    \label{fig:architecture}
    \vspace{-0.5em}
\end{figure}

This research investigates: (1)~how agents should route and validate across heterogeneous retrieval pathways---dense, sparse, and structured KG queries---under uncertainty; (2)~what provenance graph representations enable scalable entity- and span-level evidence tracking across heterogeneous sources; and (3)~whether provenance-aware KG grounding yields measurable gains in reasoning accuracy, output verifiability, and failure recovery relative to hybrid RAG baselines.

We will implement PA-AKG and develop benchmarks measuring grounded reasoning quality, evidence completeness, and failure recovery. Experiments compare vector-based RAG, hybrid retrieval, GraphRAG-style reasoning, and PA-AKG across \textbf{multi-hop QA over heterogeneous knowledge bases} and \textbf{claim verification with evidence chain reconstruction}. Expected outcomes include the PA-AKG open prototype, provenance-aware benchmarks, and baselines demonstrating that \textbf{structured provenance-aware grounding} yields measurable gains in accuracy, verifiability, and failure recovery---advancing reliable agentic AI with direct applicability to production platforms such as Amazon Bedrock Knowledge Bases.

\bibliographystyle{IEEEtran}
\bibliography{references}

\end{document}
